\subsection{Hitchin 纤维} Hitchin 的一个重要观察是：Hitchin 纤维化 $f$ 的纤维可以用「交换」的语言描述，即用 $X$ 的某些有限覆叠上的线丛来描述。下面对 $G=\GL_{n}$ 与 $G=\Sp_{2n}$ 详述这一观察。

\subsubsection{$\GL_{n}$ 的情形与谱曲线}\label{sss:spectral curve} 对 $a=(a_{1},\cdots, a_{n})\in\calA_{n,\calL}$，可以定义一条曲线 $Y_{a}$，带有 $n$ 次态射 $\pi_{a}: Y_{a}\to X$。构造如下。$\calL$ 的全空间可以写成 $X$ 上的相对谱
\begin{equation*}
\Sigma:=\Tot(\calL)=\Spec(\calO\oplus\calL^{\otimes -1}\oplus\calL^{\otimes -2}\oplus\cdots)=\Spec\Sym(\calO\oplus\calL^{\otimes -1}).
\end{equation*}
设 $\pi: \Sigma\to X$ 是投影。考虑 $X$ 上凝聚层之间的映射
\begin{equation*}
\iota_{a}: \calL^{\otimes -n}\to \pi_{*}\calO_{\Sigma}=\calO\oplus\calL^{\otimes -1}\oplus\calL^{\otimes -2}\oplus\cdots,
\end{equation*}
它在坐标下由 $((-1)^{n}a_{n}, (-1)^{n-1}a_{n-1},\cdots, -a_{1},1,0,\cdots)$ 给出。由伴随性，$\iota_{a}$ 诱导映射 $\iota'_{a}: \pi^{*}\calL^{\otimes -n}\to \calO_{\Sigma}$，其像记为 $\calI_{a}$。则 $\calI_{a}$ 是 $\Sigma$ 上的理想层。定义{\em 谱曲线} $Y_{a}$ 为 $\Sigma$ 中由 $\calI_{a}$ 定义的闭子概形
\begin{equation*}
Y_{a}=\Spec(\calO\oplus\calL^{\otimes-1}\oplus\calL^{\otimes-2}\oplus\cdots)/\calI_{a}.
\end{equation*}
若在某开子集 $U\subset X$ 上平凡化 $\calL$，并把 $a_{i}$ 看作 $U$ 上函数，则 $Y_{a}|_{U}$ 是 $U\times\AA^{1}$ 中由方程 $y^{n}-a_{1}y^{n-1}+a_{2}y^{n-2}+\cdots +(-1)^{n}a_{n}=0$ 定义的子概形（其中 $y$ 是 $\AA^{1}$ 的坐标）。投影 $\pi_{a}: Y_{a}\to X$ 是 $n$ 次有限平坦的。曲线 $Y_{a}$ 称为 $a$ 的{\em 谱曲线}，因为 $\pi_{a}$ 的纤维正是特征多项式 $y^{n}-a_{1}y^{n-1}+a_{2}y^{n-2}+\cdots +(-1)^{n}a_{n}$ 的根。

当 $a\in\Ah_{n,\calL}$ 时，曲线 $Y_{a}$ 约化，因而在某 Zariski 稠密开子集上光滑；此时存在模叠 $\cPic(Y_{a})$，分类 $Y_{a}$ 上无挠且一般秩为 $1$ 的凝聚 $\calO_{Y_{a}}$-模，见 \cite{AK}。分类 $Y_{a}$ 上线丛的通常 Picard 叠 $\Pic(Y_{a})$ 是 $\cPic(Y_{a})$ 的开子叠，它通过张量积作用在 $\cPic(Y_{a})$ 上。

\begin{prop}\label{p:cPic} 设 $a\in\Ah_{n,\calL}(k)$。设 $\calM_{a}$ 是 $f:\calM_{n,\calL}\to \calA_{n,\calL}$ 在 $a$ 上的纤维。则有叠的典范同构
\begin{equation*}
\cPic(Y_{a})\cong\calM_{a}.
\end{equation*}
\end{prop}

\subsubsection{证明梗概}\label{sss:pf cPic}
给出态射 $\cPic(Y_{a})\to \calM_{a}$。对 $Y_{a}$ 上任意凝聚层 $\calF$，正向像 $\pi_{a,*}\calF$ 是 $X$ 上凝聚层，带有映射 $\varphi_{\calF}: \pi_{a,*}\calF\otimes\calL^{-1}\to\pi_{a,*}\calF$，因为 $\calF$ 是 $\calO_{Y_{a}}$-模，而 $\calO_{Y_{a}}$ 把 $\calL^{-1}$ 作为第二个直和项包含在内。当 $\calF$ 无挠且一般秩为 $1$ 时，$\calV:=\pi_{a,*}\calF$ 在 $X$ 上无挠（因而是向量丛），秩为 $n$，且映射 $\varphi_{\calF}$ 诱导 Higgs 场 $\varphi:\calV\to\calV\otimes\calL$。对应 $\calF\mapsto (\calV,\varphi)$ 定义了态射 $\cPic(Y_{a})\to \calM_{a}$，可以证明它是同构。\qed

\subsubsection{$G=\Sp_{2n}$ 的情形} 此时 $a\in\calA_{G,\calL}$ 是多元组 $(a_{1},\cdots, a_{n})$，其中 $a_{i}\in\cohog{0}{X,\calL^{\otimes2i}}$。对 $a\in\calA_{G,\calL}$，同样可以定义谱曲线 $Y_{a}$，它是 $\calL$ 全空间中由局部生成元
\begin{equation*}
P_{a}(y)=y^{2n}+a_{1}y^{2n-2}+\cdots+a_{n}
\end{equation*}
所生成理想截出的闭子概形。注意 $Y_{a}$ 带有对合 $\sigma(y)=-y$，投影 $\pi_{a}: Y_{a}\to X$ 在此对合下不变。现设 $P_{a}(y)$ 限制到 $X$ 的一般点时可分，从而 $Y_{a}$ 约化。对合 $\sigma$ 诱导紧化 Picard 叠 $\cPic(Y_{a})$ 上的对合 $\sigma$。设 $(-)^{\vee}$ 是 $Y_{a}$ 上凝聚层的相对 Serre 对偶函子，即 $\calF^{\vee}=\un{\Hom}_{\calO_{X}}(\calF,\calO_{X})$，按自然方式看作 $\calO_{Y_{a}}$-模。设 $\cPrym(Y_{a};\sigma)$ 是数对 $(\calF,\iota)$ 的模叠，其中 $\calF\in\cPic(Y_{a})$，$\iota$ 是同构 $\iota: \sigma^{*}\calF\cong\calF^{\vee}$，且满足 $(\sigma^{*}\iota)^{\vee}=-\iota$。称 $\cPrym(Y_{a};\sigma)$ 为 $Y_{a}$ 关于对合 $\sigma$ 的{\em 紧化 Prym 叠}。与 $\GL_{n}$ 情形类似，$\calM_{a}$ 有如下描述。

\begin{prop} 设 $a\in\calA_{G,\calL}$ 使得 $Y_{a}$ 约化。则有叠的典范同构
\begin{equation*}
\cPrym(Y_{a})\cong\calM_{a}.
\end{equation*}
\end{prop}

\subsubsection{证明梗概}
对 $a\in\calA_{G,\calL}$，$\calM_{a}$ 中的点是三元组 $(\calV,\omega, \varphi)$，其中 $\calV$ 是 $X$ 上秩 $2n$ 的向量丛，$\omega: \wedge^{2}\calV\to\calO_{X}$ 是 \S\ref{sss:Bun Sp} 中 $\calV$ 上的辛形式，$\varphi:\calV\to\calV\otimes\calL$ 满足对 $\calV$ 的局部截面 $u,v$ 有 $\om(\varphi u, v)+\om(u, \varphi v)=0$，且 $\varphi$ 的特征多项式为 $P_{a}(y)$。由命题 \ref{p:cPic}，Higgs 丛 $(\calV,\varphi)$ 给出点 $\calF\in\cPic(Y_{a})$。辛形式可以看作同构 $\jmath: (\calV,\varphi)\isom (\calV^{\vee}, -\varphi^{\vee})$，满足 $\jmath^{\vee}=-\jmath$。注意 Higgs 丛 $(\calV^{\vee}, -\varphi^{\vee})$ 在命题 \ref{p:cPic} 的同构下对应 $\sigma^{*}\calF^{\vee}$。于是同构 $\jmath$ 变成 $\iota:  \sigma^{*}\calF\isom\calF^{\vee}$，满足 $(\sigma^{*}\iota)^{\vee}=-\iota$。\qed

\subsubsection{例 \ref{sss:MHilb} 续}\label{sss:Hilb}
固定 $X$ 上两个线丛 $\calL_{1}$ 与 $\calL_{2}$，设 $\calH_{\calL_{1},\calL_{2}}$ 是 \S\ref{sss:MHilb} 中模叠 $\calM_{H,\rho}$ 在 $(\calL_{1},\calL_{2})$ 上的纤维。则 $\calH_{\calL_{1},\calL_{2}}$ 分类 $(\calV,\varphi,\beta)$，其中 $\calV$ 是 $X$ 上秩 $n$ 的向量丛，$\varphi:\calU\to\calU\otimes\calL_{1}$ 是 Higgs 场，$\beta: \calU\to \calL_{2}$。有 Hitchin 型映射 $\calH_{\calL_{1},\calL_{2}}\to\calA_{n,\calL_{1}}$，把 $(\calV,\varphi,\beta)$ 送到 $(a_{i}(\varphi))_{1\leq i\leq n}$。设 $a\in\Ah_{n,\calL_{1}}$，$\calH_{\calL_{1},\calL_{2},a}$ 是 $\calH_{\calL_{1},\calL_{2}}$ 在 $a$ 上的纤维。考虑与 \S\ref{sss:spectral curve} 中相同的谱曲线 $Y_{a}$。用命题 \ref{p:cPic} 可把 $(\calU,\varphi)$ 等同于点 $\calF\in\cPic(Y_{a})$。映射 $\beta:\pi_{a,*}\calF=\calU\to\calL_{2}$ 由伴随性给出映射 $b:\calF\to \pi^{!}_{a}\calL_{2}$。这里 $\pi^{!}_{a}\calL_{2}\cong\pi^{*}_{a}\calL_{2}\otimes\omega_{Y_{a}/X}$，而相对对偶复形 $\omega_{Y_{a}/X}=\omega_{Y_{a}}\otimes\pi_{a}^{*}\omega_{X}^{-1}$ 是 $Y_{a}$ 上线丛，因为 $Y_{a}$ 是平面曲线，从而 Gorenstein。由于 $\calF$ 无挠且一般是线丛，$b$ 是凝聚层的单射。因此数据 $(\calU,\varphi,\beta)$ 变成线丛 $\pi^{!}_{a}\calL_{2}$ 的凝聚子层 $\calF$ 的数据。由于 $\pi^{!}_{a}\calL_{2}$ 是线丛，子层 $\calF$ 由商 $(\pi^{!}_{a}\calL_{2})/\calF$ 的支撑决定，而后者是 $Y_{a}$ 的零维子概形。于是得到典范同构
\begin{equation*}
\calH_{\calL_{1},\calL_{2},a}\cong\Hilb(Y_{a}),
\end{equation*}
其中 $\Hilb(Y_{a})$ 是 $Y_{a}$ 的各种长度零维子概形的 Hilbert 概形的不交并。

总之，一族平面曲线的点 Hilbert 概形可以实现在 \S\ref{sss:general Hitchin} 的框架中，作为 Hitchin 型模空间。

\subsubsection{Hitchin 纤维上的对称性}\label{sss:sym Hit}
在 $G=\GL_{n}$、$a\in\Ah_{n,\calL}$ 的情形，已把 Hitchin 纤维 $\calM_{a}$ 与谱曲线 $Y_{a}$ 的紧化 Picard 叠 $\cPic(Y_{a})$ 等同。通常 Picard 叠 $\Pic(Y_{a})$ 通过张量积作用在 $\cPic(Y_{a})$ 上。因此有 $\Pic(Y_{a})$ 在 $\calM_{a}$ 上的作用。当 $Y_{a}$ 光滑时，此作用单可迁。

在 $G=\Sp_{2n}$、$a\in\Ah_{G,\calL}$ 的情形，Prym 叠 $\Prym(Y_{a};\sigma)$ 定义为数对 $(\calL,\iota)$ 的模叠，其中 $\calL\in\Pic(Y_{a})$，$\iota:\sigma^*\calL\cong\calL^{-1}$ 满足 $(\sigma^{*}\iota)^{\vee}=\iota$。则 $\Prym(Y_{a};\sigma)$ 通过张量积作用在 $\calM_{a}\cong\cPrym(Y_{a};\sigma)$ 上。

对一般的 $G$ 与 $a\in\Ah_{G,\calL}$，同样可以定义交换群叠 $\calP_{a}$，它作用在相应的 Hitchin 纤维 $\calM_{a}$ 上。事实上，经由映射 $a:X\to [\frc/\Gm]$，\S\ref{sss:sym ASF} 中下降到 $[\frc/\Gm]$ 的正则中心化子群概形 $J$ 拉回为 $X$ 上光滑群概形 $J_{a}$，它一般地是环面。叠 $\calP_{a}$ 于是分类 $X$ 上 $J_{a}$-主齐性空间。

\subsection{与仿射 Springer 纤维的关系}\label{ss:rel Hit ASF}
本小节陈述 Ng\^o 所观察到的 Hitchin 纤维与仿射 Springer 纤维之间的精确关系。

\subsubsection{} 固定点 $a\in\Ah_{G,\calL}(k)$。把 $a$ 看作映射 $X\to [\frc/\Gm]$，设 $U_{a}\subset X$ 是 $[\frc^{\rs}/\Gm]$ 在 $a$ 下的原像。由于 $a\in\Ah_{G,\calL}$，$U_{a}$ 非空，故补集 $X-U_{a}$ 由 $X$ 的有限多个 $k$-点组成。对每个点 $x\in X-U_{a}$，设 $\calO_{x}$ 是 $X$ 在 $x$ 处的完备局部环，其分式域为 $F_{x}$，剩余域 $k(x)=k$。固定 $\calL$ 在 $x$ 附近的平凡化，可把 $a_{x}=a|_{\Spec\calO_{x}}$ 等同于 $\frc(\calO_{x})\cap\frc^{\rs}(F_{x})$ 中的元素。设 $\ep(a_{x})\in\frg^{\rs}(F_{x})$ 是它在 Kostant 截面中的对应点，则得到仿射 Grassmann 簇 $\Gr_{G, x}=L_{x}G/L^{+}_{x}G$ 中的仿射 Springer 纤维 $\sX_{a_{x}}:=\sX_{\ep(a_{x})}$（把 $x$ 放在下标中，是为了强调环路群的定义用的是域 $F_{x}$，它同构于 $k\lr{t}$）。中心化子的环路群 $LG_{\ep(a_{x})}$ 作用在 $\sX_{a_{x}}$ 上，且此作用如 \S\ref{sss:sym ASF} 那样经过局部 Picard 群 $P_{a_{x}}$ 分解。另一方面，\S\ref{sss:sym Hit} 提到整体 Picard 叠 $\calP_{a}$ 作用在 Hitchin 纤维 $\calM_{a}$ 上。Ng\^o 的乘积公式大意是说：模去局部与整体 Picard 叠的作用后，$\calM_{a}$ 是所有点 $x\in X-U_{a}$ 处仿射 Springer 纤维 $\sX_{a_{x}}$ 的乘积。

\begin{theorem}[乘积公式，Ng\^o {\cite[定理 4.6]{NgoHit} 与 \cite[命题 4.15.1]{NgoFL}}]\label{th:prod} 对 $a\in\Ah(k)$，存在典范态射
\begin{equation*}
\calP_{a}\twtimes{\prod_{x\in X-U_{a}}P_{a_{x}}}\left(\prod_{x\in X-U_{a}}\sX_{a_{x}}\right)\to \calM_{a},
\end{equation*}
它是叠的同胚。这里记号 $P\twtimes{H}Y$（其中 $H$ 右作用在 $P$ 上、左作用在 $Y$ 上）表示 $P\times Y$ 对作用 $h\cdot (p,y)=(ph^{-1}, hy)$ 的商，其中 $p\in P, y\in Y, h\in H$\zy{原文此处用 homeomorphism（同胚）一词；对象是模叠，读者可按叠意义下的同构理解，译文照原文作「同胚」。同一记号 $\twtimes{}$ 在 \S4.2 已用于 $\frc_{H,\rho,\calE_{A}}=\calE_{A}\twtimes{A}\frc_{H,\rho}$。来源：原文 \S4.4 该定理文脉与 \S4.2。}。
\end{theorem}

当 $G=\GL_{n}$ 时，乘积公式可以用谱曲线的紧化 Picard 叠以更熟悉的方式重述；后者实际上对一切约化曲线都有意义。设 $C$ 是 $k$ 上约化射影曲线。对每个点 $x\in C(k)$，可以定义局部 Picard 群 $P_{x}(C)$，其 $k$-点是 $F^{\times}_{x}/\calO^{\times}_{x}$，其中 $\calO_{x}$ 是 $C$ 在 $x$ 处的完备局部环，$F_{x}$ 是其全分式环（一般是域之积）。还有 $\cPic(C)$ 的局部类比 $\ov{P}_{x}(C)$，其 $k$-点是 $F_{x}$ 的分式 $\calO_{x}$-理想，参看 \S\ref{sss:GLn orb formula}。于是有乘积公式的如下变体，其证明与定理 \ref{th:prod} 类似。

\begin{prop}\label{p:Pic prod} 设 $C$ 是 $k$ 上约化射影曲线，$Z=C-C^{sm}$ 是 $C$ 的奇点轨迹。则有典范态射
\begin{equation*}
\Pic(C)\twtimes{\prod_{x\in Z}P_{x}(C)}\left(\prod_{x\in Z}\ov{P}_{x}(C)\right)\to\cPic(C),
\end{equation*}
它是叠的同胚。
\end{prop}

乘积公式提供了 Hitchin 纤维的几何与仿射 Springer 纤维的几何之间的联系，而后者如前所见与轨道积分密切相关。这一联系使得可以通过研究 Hitchin 纤维的上同调来接近基本引理。用 Hitchin 纤维代替仿射 Springer 纤维的好处是：Hitchin 纤维化具有更好的几何性质，这一点已在 \S\ref{sss:geom Hit} 看到。
